% 第 39 章　同时、时间与长度
\chapter{同时、时间与长度}

第 38 章把我们逼到了一个不得不接受的事实面前：真空中的光速，对一切匀速运动的观察者都是同一个数 \(c\)，与光源怎么动、观察者怎么动无关。当时我们按下了一个疑问：速度不是应该叠加吗？这一章就来付清这笔账。我们会看到，只要光速不变这一条成立，三个你从小当作常识的观念——“同时”“一秒钟”“一米”——都必须改写。这不是文字游戏，而是每天都被原子钟和宇宙射线反复验证的事实。

\step{反常识开场}

\begin{mainline}
大气层顶部，宇宙射线（主要是来自深空的高速质子）不断撞进空气分子，碎片里有一种叫 \(\mu\) 子的粒子。你可以在实验室里让 \(\mu\) 子停下来，安静地测它的寿命：平均只有 \(2.2\) 微秒，之后它就衰变成别的粒子。这是可以反复测量的实验事实。

现在做一道小学生算术题。\(\mu\) 子从大约 10 公里高空产生，就算它以光速直扑地面，\(2.2\) 微秒也只能飞
\[
3\times10^{8}\ \text{米/秒}\times 2.2\times10^{-6}\ \text{秒}\approx 660\ \text{米}\text{。}
\]
660 米，连大气层的零头都不到。按这个算法，海平面上应该几乎探测不到 \(\mu\) 子。

但事实是：每一分钟、每一平方米的地面上，都有几百个 \(\mu\) 子穿过——它们也穿过你的身体。用一只带酒精蒸气的云室，你甚至能亲眼看到它们拖出的白色径迹。粒子没有说谎，探测器也没有出错。那错的只能是我们在算术里偷偷用掉的某个假设。

第二个悬念藏在你的手机里。导航卫星上的原子钟极其精确，可工程师发现：如果不在设计上预先修正，卫星钟和地面钟每天都会对不上几十微秒。光走 1 微秒就是 300 米，几十微秒的误差累积下来，导航每天能偏出十公里量级。你的手机能把你带到正确路口，恰恰是因为有人认真对待了“天上的钟和地上的钟走得不一样快”这件事。
\end{mainline}

\step{先猜一猜}

\begin{mainline}
在继续之前，请把你的直觉押在桌上。面对“\(\mu\) 子为什么能活着到达地面”，你可能会猜：

第一，\(\mu\) 子在高速飞行时“变质”了，寿命被某种我们不懂的机制拉长了——就像剧烈运动的人心跳加速，也许粒子的“内部时钟”被速度干扰了。

第二，实验室里测的那个 \(2.2\) 微秒只适用于静止的 \(\mu\) 子，运动的 \(\mu\) 子本来就是另一种粒子。

第三，“10 公里”这个距离本身有问题——也许在 \(\mu\) 子“看来”，大气层根本没有那么厚。

还有第四种可能，也是大多数人心里默认的背景板：时间对谁都一样流，一米对谁都一样长，所以以上猜测都不靠谱，一定是实验搞错了。

请记住你的选择。本章结束时你会发现，第三种猜测离真相最近——但真相比你猜的更彻底：不是 \(\mu\) 子坏了，也不是大气层变薄了，而是“时间”和“长度”这两个词，从来就没有你以为的那种人人通用的含义。
\end{mainline}

\step{动手实验}

\begin{experiment}[看见不等于发生]
你需要一个雷雨天，或者一位在远处干活的伙伴。

\textbf{做法一}：雷雨时，看到闪电的瞬间开始默数秒数，听到雷声时停下。把秒数乘上 340 米（声音每秒大约走的距离），就是闪电到你的距离。数出 3 秒，闪电在大约 1 公里外。

\textbf{做法二}：看几十米外的人敲钉子或拍手。你会发现锤子落下的画面和敲击声对不上：画面先到，声音后到。用手机的慢动作录像回看，错位更明显。

\textbf{你在做什么}：你在区分两个时刻——“事情发生的时刻”和“我看见它的时刻”。光从闪电到你眼睛只要几微秒，可以忽略；声音却要走好几秒。所以你的大脑其实在自动执行一条规则：\emph{发生的时刻，等于看到的时刻减去信号在路上花的时间。}
\end{experiment}

\begin{mainline}
这个实验看起来和相对论毫无关系，但它藏着本章最重要的一个动作：\emph{扣除信号传播时间}。

一旦你开始认真做这件事，一连串新问题就冒出来了。假如两道闪电在你眼里同时亮起，一道来自 1 公里外，一道来自 2 公里外，它们真的是同时发生的吗？不是——扣除光在路上的时间后，远的那道其实发生得更早。反过来，两道真正同时发生的闪电，如果离你远近不同，你看到的时刻也不同。

这说明“同时”不是一个可以光靠眼睛宣布的词。要给两个异地事件排定先后，你必须规定一套操作：在每个地点放一个同步好的钟，事件在哪发生就用哪里的钟记录时刻。而“同步”两只异地的钟，又必须用信号——最忠实的信号就是光，因为它的速度最可靠（第 38 章）。爱因斯坦在 1905 年的论文开头做的正是这件看似琐碎的事：他给“同时”下了一个操作定义——用光速不变的信号来校准各地的钟。

请记住这个定义。下面所有的“怪事”，都是从“只能用光来对钟”这一个诚实的操作里长出来的。
\end{mainline}

\step{旧模型遇到麻烦}

\begin{mainline}
现在把舞台搬上一列火车。车厢正中央有一盏灯，闪烁一次，光同时向车头和车尾射去。

\textbf{车上的观察者}（相对车厢静止）这样算账：车头车尾离灯一样远，光速不变——第 38 章说过，光速与光源、观察者的运动都无关——所以两束光走相同距离、用相同时间，\emph{同时}到达两端。这不是推测：他可以在车头车尾各放一只同步好的钟，两只钟记下的到达时刻确实相同。

\textbf{地面上的观察者}看同一过程。他同样必须承认：两束光的速度都是 \(c\)，一束也不能快——这正是光速不变的内容。但就在光飞行的这段时间里，车尾在\emph{迎着}光跑，车头在\emph{逃离}光。于是光追上 车尾用的路程短，追上车头用的路程长：\emph{光先到车尾，后到车头}。他用自己的同步钟阵去记，也会记下两个不同的时刻。

\begin{center}
\begin{tikzpicture}[>=Stealth,scale=0.92]
  \draw[thick] (0,0) rectangle (8,1.3);
  \fill (4,0.65) circle (0.07) node[above=3pt]{闪光};
  \draw[->,thick,blue!60!black] (3.9,0.4) -- (1.0,0.4);
  \draw[->,thick,blue!60!black] (4.1,0.4) -- (7.0,0.4);
  \node[left] at (0,0.65){车尾};
  \node[right] at (8,0.65){车头};
  \draw[->,thick] (3,1.9) -- (5.5,1.9) node[midway,above]{火车以速度 \(v\) 向右行驶};
  \node[align=center] at (4,-0.8){地面看：车尾迎光而来，车头离光而去，\\ 两束光速度都是 \(c\)，所以先到车尾、后到车头};
\end{tikzpicture}
\end{center}

同一对事件——“光到车尾”和“光到车头”——一个参考系说同时，另一个参考系说不同时。\emph{两边都没有算错。} 注意，这和动手实验里的“看见延迟”完全是两回事：双方都已经扣除了信号传播时间，用的是摆在现场的同步钟，结论仍然不一致。

那么“到底谁对”？这个问题问不出口，因为它预设了存在一种凌驾于所有参考系之上的“真正的同时”。旧模型——从牛顿时代一路默认的“绝对时间”，写成公式就是伽利略变换里的 \(t'=t\)——在这里正式报废。不是被推翻，而是被发现：它从来只是一个近似的习惯，而不是世界的结构。
\end{mainline}

\step{发明新概念}

\begin{mainline}
“同时”的崩塌不是灾难的结束，而是重建的开始。如果时间是相对的，我们拿什么当钟？物理学家喜欢把问题推到最干净的地步：造一台原理上挑不出毛病的钟——\emph{光钟}。

光钟很简单：两面平行的镜子，相距 \(L_0\)，一束光在镜面之间来回反射，每打一个来回算“嘀嗒”一次。为什么信任它？因为光速不变，这台钟的快慢只由距离和 \(c\) 决定，没有任何零件会老化。更关键的是：如果任何别的好钟（原子钟、机械表、你的心跳）和光钟对不上，你就造出了一台能检测“自己是否在绝对运动”的仪器——这直接违背相对性原理。所以结论很强：\emph{一切}时钟，包括生物钟，都必须和光钟保持同样的走法。

先跟着钟一起飞。对你（与钟相对静止的观察者）来说，光直上直下，打一个来回的路程是 \(2L_0\)，于是一次嘀嗒是
\[
\Delta\tau=\frac{2L_0}{c}.
\]
再换到地面。现在整只光钟以速度 \(v\) 横着飞过。对你来说，光不再直上直下，它走的是一条斜线：光向上飞的这段时间里，镜子也在向前跑。设地面测得光从下到上用时 \(\Delta t/2\)，那么光走的斜边是 \(c\,\Delta t/2\)，钟在水平方向挪了 \(v\,\Delta t/2\)，而镜面间距仍是 \(L_0\)（垂直于运动方向的长度不变，下面会说明为什么）。三条边构成一个直角三角形，勾股定理给出
\[
\left(\frac{c\,\Delta t}{2}\right)^{2}=L_0^{2}+\left(\frac{v\,\Delta t}{2}\right)^{2}.
\]

\begin{center}
\begin{tikzpicture}[>=Stealth,scale=0.95]
  % 左：随钟静止的观察者
  \fill[pattern=north east lines] (-1.3,2.4) rectangle (1.3,2.56);
  \fill[pattern=north east lines] (-1.3,-0.16) rectangle (1.3,0);
  \draw[thick] (-1.3,0)--(1.3,0);
  \draw[thick] (-1.3,2.4)--(1.3,2.4);
  \draw[<->,thick,blue!60!black] (0,0.06)--(0,2.34);
  \node[right] at (0.08,1.2){$L_0$};
  \node at (0,-0.75){（甲）与钟同行：光直上直下};
  % 右：地面观察者
  \begin{scope}[xshift=7.2cm]
    \draw[thick] (-2.8,0)--(2.8,0);
    \draw[thick] (-0.7,2.4)--(0.7,2.4);
    \draw[->,thick,blue!60!black] (-1.6,0.06)--(0,2.34);
    \draw[->,thick,blue!60!black] (0,2.34)--(1.6,0.06);
    \draw[dashed] (0,0)--(0,2.4);
    \draw[dashed] (-1.6,0)--(1.6,0);
    \node[right] at (0.06,1.3){$L_0$};
    \node[below] at (-0.85,0){$\dfrac{v\,\Delta t}{2}$};
    \node[below] at (0.85,0){$\dfrac{v\,\Delta t}{2}$};
    \node[above left] at (-0.9,1.35){$\dfrac{c\,\Delta t}{2}$};
    \node at (0,-0.75){（乙）地面看：光走斜线，路程变长};
  \end{scope}
\end{tikzpicture}
\end{center}

解这个方程：把含 \(\Delta t\) 的项合并，得到
\[
\Delta t=\frac{2L_0/c}{\sqrt{1-v^{2}/c^{2}}}=\gamma\,\Delta\tau,
\qquad
\gamma=\frac{1}{\sqrt{1-v^{2}/c^{2}}}\ \ge 1.
\]
结论：在地面上看，运动光钟的每一次嘀嗒都被拉长到原来的 \(\gamma\) 倍——\emph{运动的钟走得慢}。这叫做\emph{时间膨胀}。与钟同行的观察者测出的时间 \(\Delta\tau\) 有个专门的名字：\emph{固有时间}，它是这只钟自己“亲身经历”的时间，也是所有观察者会测出的结果中最短的那个。

新名词值得停一秒。固有时间不是“真正的时间”，而是“不需要任何异地对钟就能读出的时间”——它就写在当事者自己的手腕上。这个概念会陪我们走到本章最后。

接下来是长度。请接受一个思想要求：\(\mu\) 子“能不能到达地面”是一个客观事件，所有参考系必须给出同一个答案，否则物理就乱了。在地面上我们已经算过（下一步会细算）：运动的 \(\mu\) 子寿命被拉长 \(\gamma\) 倍，所以能飞完 10 公里。现在换到\emph{跟着 \(\mu\) 子一起飞}的参考系：在这个系里，\(\mu\) 子安安静静，寿命就是老实的 \(2.2\) 微秒，只够飞 660 米——除非大气层变短了。于是唯一的出路是：在这个参考系里，迎面扑来的大气层沿运动方向\emph{收缩}了，10 公里缩成 \(10/\gamma\) 公里。这叫\emph{长度收缩}：一根静止长度为 \(L_0\)（固有长度）的尺子，在沿其长度方向以速度 \(v\) 运动的参考系中量出来是
\[
L=\frac{L_0}{\gamma}.
\]

为什么垂直方向不收缩？想象火车钻山洞：假如横向尺寸也随运动改变，车上和地面的观察者会对“火车碰不碰得到洞壁”给出不同答案——而相撞是客观事件，不许有两种答案。所以收缩只发生在运动方向上。

回头看，三个新概念——同时的相对性、时间膨胀、长度收缩——其实是同一枚硬币：它们都从“只能用光速不变的信号对钟”这一条里长出来，并且互相咬合，保证任何参考系对客观事件的判断都一致。
\end{mainline}

\step{一句话模型}

\begin{oneline}
光速对所有人相同，于是“同时”只对某个参考系成立；运动的钟按因子 \(\gamma\) 变慢，运动的尺沿运动方向按 \(\gamma\) 缩短——而固有时间和时空间隔，才是所有观察者都同意的东西。
\end{oneline}

\step{数学翻译器}

\begin{mainline}
把上面的故事压成两条公式，并按本书惯例问四个问题。
\[
\Delta t=\gamma\,\Delta\tau,\qquad
L=\frac{L_0}{\gamma},\qquad
\gamma=\frac{1}{\sqrt{1-v^{2}/c^{2}}}.
\]

\textbf{它描述什么}：同一只钟，在随它静止的系里走 \(\Delta\tau\)，在看它运动的系里走 \(\Delta t\)；同一把尺，静止时长 \(L_0\)，沿长度方向运动时量出 \(L\)。

\textbf{为什么这样写}：全部来历就是“光速不变”加勾股定理。斜边总比直角边长，所以 \(\gamma\ge 1\)：运动的钟总是变慢，运动的尺总是变短，绝不会反过来。

\textbf{现在用特殊情况做体检}：
\begin{itemize}[leftmargin=1.6em]
  \item \(v=0\)：\(\gamma=1\)，公式退回日常，理应如此。
  \item \(v\ll c\)：把根号展开，\(\gamma\approx 1+\dfrac{v^{2}}{2c^{2}}\)，偏离 1 的量是十亿分之一量级往下——这就是为什么人类三百年都没发现牛顿力学有缝。
  \item 方向反转：公式里只有 \(v^{2}\)，所以钟迎面飞来和离你而去，变慢一样多。时间膨胀跟“看起来快还是慢”的多普勒效应是两码事。
  \item \(v\to c\)：\(\gamma\to\infty\)，运动的钟趋于停摆。光本身的“钟”不流逝固有时间——这句话点到为止，它的严格含义要等时空间隔。
  \item \(v>c\)：根号里变成负数，公式拒绝给出实数答案。这不是数学的坏脾气，而是在说：普通物体到不了光速。为什么到不了，是第 40 章能量与动量的故事。
\end{itemize}

\textbf{何时成立、何时失效}：惯性参考系、平直时空内严格成立；涉及加速度和引力时要升级（第 41 章）。下面算几笔真实的账。

\textbf{账单一：汽车。} 时速 120 公里约 33 米/秒，
\[
\gamma-1\approx\frac{v^{2}}{2c^{2}}\approx\frac{33^{2}}{2\times(3\times10^{8})^{2}}\approx 6\times10^{-15}.
\]
你开一辈子车，“赚”到的时间不过纳秒量级。直觉没错，只是适用范围太窄。

\textbf{账单二：\(\mu\) 子。} 高空产生的 \(\mu\) 子速度约 \(0.998c\)：
\[
\gamma=\frac{1}{\sqrt{1-0.998^{2}}}\approx 16.
\]
地面看：寿命 \(2.2\times16\approx 35\) 微秒，飞行距离 \(0.998\times300\ \text{米/微秒}\times35\ \text{微秒}\approx 10.4\) 公里——刚好够到海平面。

\textbf{账单三：导航卫星。} 卫星速度约 3.9 公里/秒，\(\gamma-1\approx 8\times10^{-11}\)，每天累积约 7 微秒的变慢。完整的修正里还有一项引力效应（第 41 章）让星载钟每天变快约 45 微秒，合计每天快约 38 微秒。工程师的做法很朴素：卫星上天前，把钟的振荡频率按理论值预先调偏。你手机里的每一次定位，都是相对论发的工资。

\textbf{账单四：环球飞行的原子钟。} 1971 年，哈菲勒和基廷把原子钟带上民航客机绕地球飞，回来和地面同款钟对比，东西两个方向的时差不同，且都在几十到几百纳秒量级上与理论（含引力项）吻合。时间膨胀从思想实验变成了机票钱能买到的实验。
\end{mainline}

\begin{mainline}
最后一件数学装备，回答“到底有什么是不变的”。定义\emph{时空间隔}
\[
s^{2}=c^{2}\Delta t^{2}-\Delta x^{2}.
\]
如果两个事件由一束光连接，那么 \(\Delta x=c\,\Delta t\)，于是 \(s^{2}=0\)——而光速不变保证了这件事在所有参考系都成立。更进一步可以证明（数学桥层）：对任意两个事件，\(s^{2}\) 在所有惯性系中都相同。当 \(s^{2}>0\) 时，\(\sqrt{s^{2}}/c\) 正是沿直线从一事到另一事的观察者测到的固有时间。旧世界里人人共享的“时间”和“长度”退位了，新的绝对量是这个带减号的组合。
\end{mainline}

\begin{mathbridge}
\textbf{洛伦兹变换。} 设两个惯性系的坐标轴平行，\(K'\) 相对 \(K\) 以速度 \(v\) 沿 \(x\) 方向运动，则同一事件的两套坐标满足
\[
x'=\gamma\,(x-vt),\qquad
t'=\gamma\left(t-\frac{vx}{c^{2}}\right).
\]
两式立刻偿还本章的三笔账。第一，同时性的相对性：在 \(K\) 中同时（\(\Delta t=0\)）但异地（\(\Delta x\neq 0\)）的事件，在 \(K'\) 中 \(\Delta t'=-\gamma v\,\Delta x/c^{2}\neq 0\)。第二，对变换直接计算可得 \(c^{2}\Delta t'^{2}-\Delta x'^{2}=c^{2}\Delta t^{2}-\Delta x^{2}\)，即间隔不变。第三，时间膨胀与长度收缩都是它的两行推论。另外，当 \(v\ll c\) 时 \(\gamma\to 1\)、\(vx/c^{2}\to 0\)，变换退化为伽利略变换 \(x'=x-vt\)、\(t'=t\)——旧模型是新模型在低速下的极限，而不是对手。
\end{mathbridge}

\step{模型的边界}

\begin{mainline}
这套模型漂亮，但有明确的门槛和几处经典的误读区。

\textbf{适用边界。} 本章一切结论的前提是惯性参考系和平直时空。一旦牵涉持续的加速、转向，或者强引力场，就要换更大的框架（第 41 章）。在工程上：时速几百公里以内，修正量在 \(10^{-15}\) 量级，盖楼、造车、设计桥梁尽可用牛顿；但导航卫星和粒子加速器不行——加速器里电子的 \(\gamma\) 可以到几万，按牛顿力学设计的磁铁根本兜不住束流。

\textbf{误读一：“一切都是相对的”。} 恰恰相反。本章的主角其实是\emph{不变量}：光速不变、固有时间不变、时空间隔不变、类时事件的先后次序不变。相对的是那些我们误以为绝对的旧概念。称它为“相对论”是历史原因，叫它“不变量理论”反而更贴切。

\textbf{误读二：“钟慢到底谁慢”。} 甲看乙的钟慢，乙看甲的钟也慢，两人都正确——因为比较两只异地的钟必须借助“同时”，而两人对“同时”的判定不同。这没有矛盾，就像两个人都说对方在自己的右边。要让分歧变成事实，两人必须重逢当面比钟；而重逢要求至少一人掉头，对称性就此打破。这就是下面双生子故事的钥匙。

\textbf{误读三：“长度收缩是看起来压扁了”。} 收缩是用同步钟阵\emph{测量}的结果，不是照片。照片记录的是同时到达底片的光子，这些光子并非同时从物体发出。仔细计算会发现，高速运动的立方体拍出来并不扁，反而像转过了一个角度（特瑞尔转动）。你“看见”的和参考系“测量”的，永远是两笔账。

\textbf{误读四：“钟被加速度弄坏了”。} 时间膨胀不是钟的机械故障。它对原子钟、心跳、化学反应一视同仁，因为它描述的是时间本身的结构。一只“被弄坏”的钟反而会立刻露馅：它会和别的钟对不上，而我们从未见过这种露馅。

最后声明一个表述约定：本书全程使用\emph{不变质量} \(m\)，不说“质量随速度变大”。高速物体真正增长的是能量和动量，那是第 40 章的主题。
\end{mainline}

\step{回头解释开场现象}

\begin{mainline}
现在把开场的两笔账结清，而且用两种参考系各结一次——它们必须一致。

\textbf{\(\mu\) 子，地面视角}：高速 \(\mu\) 子的钟变慢，寿命从 \(2.2\) 微秒膨胀到约 \(35\) 微秒，以 \(0.998c\) 飞行约 10.4 公里，到达海平面没有问题。

\textbf{\(\mu\) 子，自己的视角}：我的寿命就是老实的 \(2.2\) 微秒，一点没变。但迎面扑来的大气层以 \(0.998c\) 运动，10.4 公里的厚度收缩成 \(10.4/16\approx 0.66\) 公里。这段路以 \(0.998c\) 扫过我只要约 \(2.2\) 微秒——刚好来得及。

两种说法一个谈时间，一个谈长度，看似毫不相干，给出的却是同一个结局：\(\mu\) 子到达地面。这不是巧合，而是时空间隔不变在背后兜底。\(\mu\) 子从产生到衰变，自己的固有时间永远是 \(2.2\) 微秒，任何参考系算出来都一样。

\textbf{导航卫星}：工程师不等卫星出错再修，而是把本章（和第 41 章引力部分）的修正直接写进设计：星载钟出厂时的频率被预先调偏约十亿分之四，上天之后恰好与地面同步。你手机地图上那个稳定的小蓝点，是“运动的钟走得慢”这条反常识结论的日常证人。

回到“先猜一猜”。猜“\(\mu\) 子变质了”的读者方向不对：粒子无辜，它只是在按自己的时间活着。猜“大气层没那么厚”的读者摸到了门把手——但那扇门后面不是错觉，而是时空本身的新结构。
\end{mainline}

\step{脑洞实验与挑战题}

\begin{thought}[双生子：谁年轻，凭什么]
甲留在地球，乙乘飞船以 \(0.8c\)（此时 \(\gamma=5/3\)）飞往 4 光年外的恒星，到达后立即掉头以同样速度返回。

地球账本：单程 \(4/0.8=5\) 年，往返共 10 年。飞船账本：飞船上固有时间缩短 \(\gamma\) 倍，乙全程只经历 \(10\times 3/5=6\) 年。重逢时，乙比甲年轻 4 岁。两人的生物钟、皱纹、记忆都同意这个结果——这是面对面比出来的，不靠任何信号约定。

可是且慢：分离期间乙不也看甲的钟慢吗？乙为什么不觉得甲应该更年轻？钥匙在“掉头”。甲全程呆在一个惯性系里；乙在掉头那一刻必须减速、反向、再加速，从一个惯性系跳到另一个惯性系。两人的经历不对称，结论当然允许不对称。分离阶段“互看对方钟慢”依然句句属实，但那段比较依赖各自的“同时”约定，而重逢时的事实不依赖任何约定。
\end{thought}

\begin{challenge}
\textbf{直线路径的固有时最长。} 时空间隔给双生子一个干净的几何表述：在时空图上画出起点（出发）和终点（重逢），甲走的是一条“直线”（惯性运动），乙走的是折线。可以证明：连接两个类时分离事件的所有路径中，惯性路径的固有时间\emph{最长}——任何绕路（哪怕只是一小段加速再折返）都会让亲历者老得更少。注意这与日常几何正好相反（欧氏几何里直线最短），负号全藏在 \(s^{2}=c^{2}\Delta t^{2}-\Delta x^{2}\) 那个减号里。闵可夫斯基几何与欧氏几何的全部“反常”，都从这一个小小的减号开始。
\end{challenge}

\begin{mainline}
\textbf{因果的边界：光锥。} 最后一个概念把本章收拢成一张图。以“此时此地”为原点，横轴是空间 \(x\)，纵轴是时间乘以 \(c\)。从原点出发的光走 \(x=\pm ct\)，在图上是两条斜线，把平面分成三个区域。上方区域里的点满足 \(c^{2}t^{2}>x^{2}\)，你以低于光速的速度来得及赶到——那是你的\emph{未来}；下方是你可能来自的\emph{过去}；左右两侧是 \(s^{2}<0\) 的\emph{类空}区域，那里的点与你之间的间隔连光都跑不完，此刻的你既影响不了它，它（在此刻）也影响不了你。

\begin{center}
\begin{tikzpicture}[>=Stealth,scale=1.25]
  \fill[blue!12] (0,0) -- (1.6,1.6) -- (-1.6,1.6) -- cycle;
  \fill[orange!15] (0,0) -- (1.6,-1.6) -- (-1.6,-1.6) -- cycle;
  \draw[->] (-2.2,0) -- (2.2,0) node[right] {$x$};
  \draw[->] (0,-1.85) -- (0,2.0) node[above] {$ct$};
  \draw[thick] (-1.6,-1.6) -- (1.6,1.6);
  \draw[thick] (-1.6,1.6) -- (1.6,-1.6);
  \node at (0,1.05){未来};
  \node at (0,-1.0){过去};
  \node at (1.45,0.3){类空};
  \node at (-1.45,0.3){类空};
  \fill (0,0) circle (0.045) node[below right=1pt]{此时此地};
\end{tikzpicture}
\end{center}

这张图给出一条比“光速不能超越”更深的表述：\emph{光锥是因果关系的边界}。对类时分开的两个事件，所有参考系都同意谁前谁后，所以“原因先于结果”是绝对的；对类空分开的两个事件，先后次序可以随参考系翻转——但这不会造成因果混乱，因为它们之间本来就不可能传递任何影响。假如存在超光速信号，某些参考系里“回信”会先于“提问”到达，因果律当场崩塌。光速上限守的不是光的特权，而是因果本身。
\end{mainline}

\begin{mainline}
\textbf{一个极端尺度的例子。} 银河系直径约 10 万光年。设想一艘 \(\gamma=10^{5}\) 的飞船（速度只比光速小约百亿分之五）横穿银河：地球看这趟旅程要 10 万年，船员却只经历大约 1 年——沿途星光在他们前方被压进极高的频率，航程的真实代价是天文数字的能量（第 40 章会算这笔账）。时间和长度的相对性不是小数点后十几位的学问，在足够大的 \(\gamma\) 下，它把“10 万年”和“1 年”变成同一段旅程的两种读数。

本章连接后文的方式也值得一提：第 40 章将从间隔不变量推出相对论的能量与动量，并回答“为什么加速到光速要付出无穷大的能量”；第 41 章会问，如果引力也能让钟变快变慢（GPS 里那项每天 45 微秒的修正），“惯性系”这个概念本身还站得住吗。而“固有时间”这个看似朴素的概念，会在量子篇以另一种身份回来。
\end{mainline}

\begin{puzzles}
  \item 火车与隧道。一列固有长度 100 米的火车以 \(0.8c\) 穿过一座固有长度 60 米的隧道。地面管理员说：“火车收缩到 60 米，有一瞬间整列火车都在隧道里。”车上乘客说：“隧道收缩到 36 米，火车什么时候都不可能全进去。”两人都错了吗？都没有吗？\emph{提示}：检查“整列火车都在隧道内”这句话——它要求车头未出、车尾已进，两个事件在两地发生，想想“同时”二字在两个人各自的参考系里意味着什么。

  \item 事件 A 发生在原点、时刻为零；事件 B 在 \(K\) 系中位于 3 光秒外、2 秒之后。是否存在某个惯性系，其中 B 早于 A？若 B 改为 1 光秒外、2 秒之后呢？\emph{提示}：先算两种情形的时空间隔 \(s^{2}\)，再想想光锥图里哪些区域允许先后次序随参考系翻转。

  \item 两艘飞船擦肩而过、互相远离，甲用望远镜看乙的钟，发现它慢得更多，比公式预言的还多。甲是否推翻了时间膨胀？\emph{提示}：望远镜里“看到”的时刻混着信号传播时间，分开的光除了钟慢，还有什么效应在排队？想想多普勒。

  \item 一台固有长度 1 米的立方体以接近光速从你面前掠过，你按下快门。照片上的它是什么形状？\emph{提示}：快门记录的是同一时刻\emph{到达}底片的光子。远端和近端的光子不是同时发出的，把这段路程差和收缩叠在一起再看。
\end{puzzles}
